\documentclass[12pt]{article}
\usepackage[utf8]{inputenc}
\usepackage[T1]{fontenc}
\usepackage[margin=1in]{geometry}
\usepackage{setspace}
\usepackage{xcolor}
\usepackage[colorlinks=true, urlcolor=blue!60!black, linkcolor=blue!60!black]{hyperref}
\setlength{\parindent}{2em}
\setlength{\parskip}{1.25em}
\renewcommand{\baselinestretch}{1.0}
\begin{document}
\begin{center}
{\LARGE \textit{Trump v.\ Cook}, 609 U.S.\ 528 (2026)\par}
\vspace{0.8em}
{\large Case Brief --- Four-Page Edition\par}
\end{center}
\vspace{1em}

\section{Memory Jogger}

On the day it overruled \textit{Humphrey's Executor}, the Court, 5--4 on the interim docket, upheld Federal Reserve Governors' ``for cause'' removal protection and kept Governor Lisa Cook in office because the President gave her no pretermination notice or chance to respond under 12 U.S.C.\ \S~242 (\href{https://supreme.justia.com/cases/federal/us/609/25a312/}{Justia}).

\section{Detailed Case Facts}

Under \S~242, each Federal Reserve Governor ``shall hold office for a term of fourteen years from the expiration of the term of his predecessor, unless sooner removed for cause by the President.'' The Fed funds itself outside appropriations and sets monetary policy through the Federal Open Market Committee. President Biden appointed Lisa Cook in 2022, and she was confirmed in 2023 to a term ending in 2038. On August 20, 2025, FHFA Director William Pulte posted a referral letter accusing Cook of 2021 mortgage fraud---declaring both a Michigan and a Georgia property her principal residence two weeks apart. Within 30 minutes President Trump posted, ``Cook must resign, now!!!'' On August 25 he sent a letter removing her ``for cause'' effective immediately, citing ``reason to believe'' she made false statements and showing at least ``gross negligence.'' Cook was the first Governor a President had tried to fire in the Fed's 111-year history. She denied wrongdoing and received no notice of the evidence, no chance to respond, and no hearing.

\section{Procedural History}

Cook sued the President, Chair Jerome Powell, and the Board. Judge Jia Cobb granted a preliminary injunction, holding that ``cause'' reaches only in-office conduct and that Cook had a due process property interest in her office. A divided D.C.\ Circuit denied a stay, Judge Katsas dissenting. The Government sought an emergency stay; the Court deferred action, heard argument on January 21, 2026, and denied the stay on June 29, 2026.

\section{Judicial Votes}

\textbf{5--4}, stay denied.
\begin{itemize}
\item \textbf{Majority:} Chief Justice Roberts, joined by Justices Sotomayor, Kagan, Kavanaugh, and Jackson.
\item \textbf{Concurrences:} Justice Kavanaugh; Justice Jackson.
\item \textbf{Dissents:} Justice Thomas; Justice Alito, joined by Justice Gorsuch; Justice Barrett.
\end{itemize}

\section{Holding}

The Government is unlikely to succeed because the President did not provide the pretermination notice and opportunity to respond that \S~242 requires, so the removal letter was void; the Court did not reach due process. Along the way it held that ``for cause'' removals are judicially reviewable; that ``cause'' is a ``substantial threshold'' turning on the seriousness of the misconduct and its nexus to official duties, not a pretext for a ``more congenial'' replacement; that courts may enter interim injunctions keeping an incumbent in office; and that the Governors' protection is consistent with Article~II. It did not decide whether the mortgage allegations, if proven, would be cause.

\section{Analysis of Opinions}

\textbf{Majority (Roberts).} Tracing central-bank independence from the First and Second Banks to the Fed, the Court stopped at likelihood of success. Under \textit{Loper Bright} and the common law, the boundary of ``cause'' is for courts. The Government's low bar would invite pretextual firings; Cook's in-office-only bar lacked textual basis. Equity can protect a \textit{de facto} incumbent pending litigation, distinguishing \textit{In re Sawyer}. Congress enacted \S~242 against \textit{Reagan} and \textit{Shurtleff}, under which ``notice and hearing are essential'' for fixed-term, for-cause offices; a social-media post fell far short. The Fed's protection preserves a founding-era balance without being ``trapped in amber.''

\textbf{Kavanaugh (concurring).} The President may still try to remove Cook with proper process, and leaving the Fed's status open after \textit{Slaughter} would unsettle markets.

\textbf{Jackson (concurring).} The equities alone defeat the stay; keeping an official in office is not irreparable injury to the President.

\textbf{Thomas (dissenting).} Office is not property; mortgage fraud is plainly cause; the notice rule rests on a ``single line of dicta''; the Fed is a modern executive agency whose protection is unconstitutional; and courts lack power to enjoin removals.

\textbf{Alito, joined by Gorsuch (dissenting).} The Court decided too much on a thin record; it should have reversed the lower courts' actual rationales and granted the stay.

\textbf{Barrett (dissenting).} Validating the Fed was ``entirely outside the scope of this case'' and ``in serious tension with \textit{Trump v.\ Slaughter},'' and the injunction is overbroad.

\section{Examples: Future Applications}

\begin{enumerate}
\item \textbf{Same-Side: Firing Over Rate Votes.} A Governor fired overnight for ``gross economic mismanagement'' after refusing to cut rates gets an injunction: no process, and a pretextual policy disagreement.
\item \textbf{Same-Side: Trivial Pre-Office Lapse.} Even with notice, firing a Governor for a late \$150 property-tax filing years earlier lacks seriousness and nexus.
\item \textbf{Opposite-Side: Insider Trading.} A Governor who pleads guilty to leaking FOMC decisions, after written notice and a chance to respond, may be removed.
\item \textbf{Opposite-Side: SEC Commissioner.} Officials without the Fed's lineage are removable at will under \textit{Slaughter}.
\item \textbf{Fence-Sitter: Renewed Mortgage Case.} With full written process and a rebuttal claiming inadvertence, a finding of ``gross negligence'' forces the open questions of deference and whether pre-office negligence suffices.
\end{enumerate}

\section{Critique}

\textbf{Incoherence with \textit{Slaughter}.} Justice Thomas argued that the First and Second Banks exercised no sovereign power, making the analogy to the modern Board ahistorical; textualists faulted reading a hearing requirement into \S~242; and Alito, Gorsuch, and Barrett objected to deciding a waived constitutional question on the emergency docket.

\textbf{Unanswered questions.} Progressive commentators welcomed the result but asked why the logic of independence does not protect other regulators. Daniel Tarullo and the Congressional Research Service noted uncertainty about review after a renewed removal and footnote 7's warning about powers ``attenuated from monetary policy.''

\section{Key Quotations}

\begin{quote}
``At minimum, Cook was entitled to some explanation of the evidence at issue, some avenue for a response, and a deadline by which a response would be due.''
\end{quote}
\noindent --- Chief Justice Roberts, for the Court, slip op.\ at 24.

\begin{quote}
``Not only the fact of independence but also the appearance of independence is key to the Federal Reserve's design. \ldots{} Without such constraints in place, any perceived or alleged misstep (past or present) could provide a ready pretext for a Governor's removal---a fact that he would surely know, and that would surely weigh on him as he decided what to say and how to vote. Nothing could be more corrosive of the independence that Congress sought to preserve.''
\end{quote}
\noindent --- Chief Justice Roberts, for the Court, slip op.\ at 14--15.

\begin{quote}
``To accept any one of those arguments would in effect transform the Federal Reserve's for-cause protection into at-will employment---an interpretive leap out of step with the statute Congress enacted and our Nation's tradition of central banking protected from political interference.''
\end{quote}
\noindent --- Chief Justice Roberts, for the Court, slip op.\ at 9.

\end{document}
